\section{Limitations and Future Work}\label{sec-9}

We surface six concrete limitations and three deliberate scope exclusions
of v1. Each is paired with the section that pre-empts the related
anticipated objection.

\textbf{(L1) MAI-DxO \ensuremath{\times} GPT-5 and DeepRare \ensuremath{\times} GPT-5 cells missing.} Two
(agent \ensuremath{\times} backbone) cells are systematically incompatible. (a) MAI-DxO's
panel orchestration over 3 deliberation iterations exceeded our 600 s
subprocess cap on GPT-5 in every pilot case, even with
\texttt{reasoning\_effort=minimal} propagated. (b) DeepRare's local-embedding
pipeline calls \texttt{eval\_tokenizer(diseases,\ max\_length=36)} on the
agent-emitted candidate list; GPT-5 at \texttt{reasoning\_effort=minimal}
emits an empty \texttt{diseases} list with high frequency, raising
\texttt{IndexError} in \texttt{transformers.tokenization\_utils\_fast} (50/50 cases
failed in Phase 2 v2). Both failures share a broader pattern ---
\textbf{frontier reasoning models with reasoning forced off may under-emit content in agent scaffolding loops that consume model output downstream}. We exclude these two cells rather than the full GPT-5
row; the same cells on DeepSeek V3.2 and Gemini 3 Flash complete in
90--180 s (MAI-DxO) and 25/25 ok (DeepRare). The pattern is
reasoning-budget \ensuremath{\times} downstream-consumer specific and tractable with
\texttt{reasoning\_effort=high}, but that re-introduces the silent
\texttt{max\_tokens} consumption we documented in \Cref{sec-5-2}.

\textbf{(L2) PhenoBrain dropped from agent lineup.} PhenoBrain (Sun et al.,
\emph{Nat. Commun.} 2025) was on our original scout list. Its 14 GB Google
Drive checkpoint is hosted under nested sub-folders whose permission
chain blocks programmatic \texttt{gdown} listing; we obtained 7.6 GB of the
checkpoint by manual download before deciding the partial state was
unsafe to evaluate. We replaced PhenoBrain with LIRICAL (the canonical
classical Bayesian baseline) in the v1 lineup. Future work re-includes
PhenoBrain once the upstream authors confirm a public mirror.

\textbf{(L3) Chinese diagnostic layer not evaluable from public data.} RareBench's
PUMCH-ADM subset is drawn from a 1,650-case single-center cohort at Peking
Union Medical College Hospital (PUMCH; 1,183 rare + 467 common). We
investigated whether a Chinese layer could be built and found two compounding
barriers. First, the labelled corpus (cases carrying the \texttt{RareDisease}
OMIM/Orphanet/CCRD ground truth) is not publicly distributed; the source
publication provides no data-access application channel or data-use agreement, reporting only an ethics approval (S-K2051), so obtaining it would
require direct institutional collaboration rather than a self-serve download.
Second, the \emph{only} PUMCH data that is public --- the public 87-case PUMCH slice
released with the RareBench code --- contains phenotype annotations only (Chinese entity strings paired one-to-one with standardised HPO terms)
and no diagnosis labels, so it cannot be scored for differential
diagnosis. A Chinese diagnostic layer (and hypothesis H5 on English-anchoring
bias) is therefore deferred to v2, contingent on the restricted labelled
corpus. The Chinese-language slot is already wired through our canonical case
model and data loader to enable a one-call extension once such data is obtained
(Future Work item 2).

\textbf{(L4) Pillar 4 (family-aware) and H9 not testable on the current corpus.} We verified (2026-05-29) that our ingested data carries no
structured family/pedigree or inheritance-mode signal: the
Phenopacket-Store cohort export we use has 0/200 files with a pedigree
block, and the family-structure and inheritance-mode fields are
unpopulated across all four layers.
Independently, no agent in our lineup performs \emph{family-aware diagnosis}
(DeepRare and MAI-DxO explicitly do not consume pedigrees; only RDMA
exposes a family-history toggle, and it is a phenotype-extraction, not a
diagnosis, component). Pre-registered H9 (``family-aware gains accrue
only on AR/XL cases'') therefore cannot be evaluated without a new
pedigree-bearing corpus (e.g.~MyGene2 / DDD) and a Pillar-4 diagnosis
path; both are deferred to v2. v1 reports four pillars (P1--P3 + P5).

\textbf{(L5) HPO extraction silver gold is LLM-generated.} Our Pillar 1
end-to-end evaluation (\Cref{sec-7-1}) compares extractors against an Opus 4.7
silver-gold reference rather than physician-annotated gold. We mitigate
this in three ways: (a) Opus is held out from the agent and judge
lineup so no agent is judged against its own backbone; (b) we report
inter-rater Jaccard between Opus and Gemini 3 Flash silver-gold (0.41)
to surface non-redundancy of the reference; (c) the 200-case PMC OA
post-cutoff holdout is being annotated by the user (clinical co-author)
to convert to physician-validated gold for the camera-ready revision.
As an interim stand-in (Ablation A5, \Cref{sec-8}), a frontier-model agent (Claude Opus 4.8) independently verified all 198 held-out cases against
the full paper text: it concurred with the Gemini extractor's \emph{diagnosis}
in 198/198 cases and judged the extracted HPO 90.4 \% precise
(8.7 \% of terms flagged unsupported --- mostly negation and subtype errors),
with a recall gap (687 salient phenotypes it judged missed). At
camera-ready we recompute A5 with the physician gold and report
Opus-vs-physician Cohen's \ensuremath{\kappa}.

\textbf{(L6) Cost reporting heterogeneity.} Six of eight adapters route
through our OpenRouter wrapper and report exact USD per call. Three adapters (RDMA, LIRICAL,
VC-RDAgent) call backbones outside the wrapper; we estimate their cost
from logged token counts and the OpenRouter price table. The
estimation introduces a \ensuremath{\leq}5 \% error band on the cost-vs-accuracy
scatter (\Cref{fig:fig2_cost_accuracy}) and we annotate estimated cells with \texttt{\dag{}} in
Appendix J. Full per-cell cost analysis is in Appendix J.

\textbf{(L7) Bounded data-contamination signal on LLM backbones.} Our A6
TS-Guessing audit (\Cref{sec-7-10} / \Cref{sec-8-9}) finds Spearman \ensuremath{\rho} between log
pre-cutoff PubMed mention count and per-disease R@1 of 0.29--0.37
across all four LLM backbones (Gemini 3 Flash, GPT-5-minimal, DeepSeek
V4-Flash, V4-Pro), versus \ensuremath{\rho} \ensuremath{\approx} 0 on classical/offline baselines
(LIRICAL, VC-RDAgent --- methodological control). This means LLM R@1 is
\emph{weakly} but consistently elevated on diseases that were better
represented in pre-cutoff literature: a real but bounded
training-frequency effect explaining \ensuremath{\approx} 9 \% of R@1 variance (\ensuremath{\rho}\textsuperscript{2} \ensuremath{\approx} 0.09).
The pre-cutoff layers (L1--L3) headline numbers therefore carry a
\ensuremath{\leq}10 \% residual contamination band. A second, independent test (H3, \Cref{sec-7-10-1}) directly compares a difficulty-matched pre- vs post-cutoff PMC set built with the identical pipeline (same source, query,
extractor, gold): pooled Gemini R@1 is 0.57 pre-cutoff vs 0.62 post-cutoff --- performance does \emph{not} drop on genuinely unseen 2024+
cases (if anything rises), so memorisation is not the driver. We make the
correlation transparent rather than dropping the pre-cutoff layers,
because (a) classical baseline \ensuremath{\rho} \ensuremath{\approx} 0 supplies a tight upper bound on
how much memorisation could explain (b) F1 (classical \textgreater{} LLM on the
rarest tier) is in the \emph{opposite} direction of training exposure and
therefore not driven by it.

\subsubsection{\texorpdfstring{Deliberate scope exclusions (these are \emph{not} defects --- they are scope choices)}{Deliberate scope exclusions (these are not defects --- they are scope choices)}}\label{deliberate-scope-exclusions-these-are-not-defects-they-are-scope-choices}

\textbf{(S1) Retrospective, not prospective.} We frame the benchmark as
\textbf{retrospective decision support} evaluation, not autonomous
diagnosis. No clinical claims are made. Anticipated objection \#12 is pre-empted in \Cref{sec-10} Conclusion.

\textbf{(S2) No fine-tuned medical LLM baseline.} Meditron-70B and
OpenBioLLM-70B are not in the v1 backbone set --- we focus on frontier
general-purpose backbones plus DeepSeek V3.2 as the open-weight
contrast. Adding a fine-tuned medical baseline is Future Work item 3;
the harness is backbone-agnostic and the addition is purely a
deployment task.

\textbf{(S3) Single-image / multimodal pillar absent.} Rare-disease
phenotyping increasingly involves facial photograph, MRI, or
electron-microscopy signal. v1 is text-only. v2 adds a multimodal
pillar (Future Work item 4).

\subsubsection{Future work (in order of expected impact)}\label{future-work-in-order-of-expected-impact}

\begin{enumerate}
\def\labelenumi{\arabic{enumi}.}
\tightlist
\item
  Pillar 4 (family-aware) --- once Phenopacket-Store pedigree fields
  are normalised + MyGene2 / DDD access converges, \textasciitilde1 month effort.
\item
  Chinese rare-disease layer (PUMCH-ADM) --- once institutional
  agreement returns, drop-in dataset add.
\item
  Fine-tuned medical LLM backbone --- add Meditron-70B and
  OpenBioLLM-70B to the backbone-axis grid (A8 extension).
\item
  Multimodal pillar (P6) --- facial / radiograph inputs; requires
  new agent set (face2gene, \ldots).
\item
  Prospective clinical study --- partner with rare-disease referral
  centre to A/B-test scaffolded agents against current-of-care
  workup. This is the long-term clinical-deployment ask the v1
  benchmark is designed to \emph{enable}, not \emph{substitute}.
\end{enumerate}

\subsubsection{Cross-references to pre-empted anticipated objections}\label{cross-references-to-pre-empted-anticipated-objections}

\begin{table*}[t]
\centering\scriptsize\setlength{\tabcolsep}{3.5pt}
\caption{Anticipated objections and the section that pre-empts each.}\label{tbl:objections}
\begin{tabular}{@{}
>{\raggedright\arraybackslash}p{(\linewidth - 2\tabcolsep) * \real{0.5000}}
  >{\raggedright\arraybackslash}p{(\linewidth - 2\tabcolsep) * \real{0.5000}}@{}}
\toprule
\begin{minipage}[b]{\linewidth}\raggedright
Attack
\end{minipage} & \begin{minipage}[b]{\linewidth}\raggedright
Where addressed
\end{minipage} \\
\midrule
\#1 Data contamination & \Cref{sec-7-10} + \Cref{sec-8-9} (A6) + \Cref{sec-9} L7 + Appendix D \\
\#2 Heterogeneous-agent fairness & \Cref{sec-5-1} Agent Fairness Matrix \\
\#4 Statistical rigor & \Cref{sec-6} footnote (Holm--Bonferroni + bootstrap CI) \\
\#5 MIMIC ICU bias & \Cref{sec-4-2} four-layer stack \\
\#7 Arbitrary agent selection & \Cref{sec-5-4} pre-registration of inclusion criteria \\
\#8 Multi-agent doesn't always help & \Cref{sec-7-2} + headline finding F2 \\
\#9 Cost not clinically meaningful & \Cref{sec-6-3} three-axis cost reporting \\
\#10 LLM-judge unreliable & \Cref{sec-7-5} + Ablation A12 \\
\#11 Model version changes silently & \Cref{sec-5-2} dated aliases + Appendix N Docker hash \\
\#3, \#6, \#12 & \Cref{sec-9} above (L2, L3, S1) \\
\bottomrule
\end{tabular}
\end{table*}

